% BEGIN chapters/introduction.tex
\chapter*{导言}
\addcontentsline{toc}{chapter}{导言}
\markboth{导言}{导言}
\src{5—9}
\begin{flushright}
献给 Oscar Zariski 与 André Weil。
\end{flushright}

本篇论文以及计划陆续发表的许多后续论文，旨在构成一部关于代数几何基础的专著。
原则上，我们不预设读者对这一学科有任何专门知识。事实甚至表明，虽然这类知识显然有其益处，
但有时也可能妨碍读者熟悉本书所阐述的观点与技术，因为它容易使人养成过于偏重双有理观点的习惯。
另一方面，我们将假定读者对下列主题具有良好的知识基础：
\begin{enumerate}[label=\alph*)]
\item 交换代数。例如 N. Bourbaki 的《数学原本》中正在编写的有关各卷所阐述的内容；
在这些卷册出版以前，可参阅 Samuel—Zariski\cite{b13} 以及 Samuel\cite{b11,b12}。
\item 同调代数。我们参照 Cartan—Eilenberg\cite{b2}（记作 (M)）、Godement\cite{b4}（记作 (G)），
以及 A. Grothendieck 最近的论文\cite{b6}（记作 (T)）。
\item 层论。主要参考文献仍为 (G) 与 (T)。这一理论提供了不可或缺的语言，使我们能够用“几何”术语
解释交换代数的基本概念，并将这些概念“整体化”。
\item 最后，读者若对函子语言有所熟悉，将会很有帮助；本专著将不断使用这种语言。
有关内容可参阅 (M)、(G)，尤其是 (T)。这一语言的原理以及函子一般理论的主要结果，
将在本专著两位作者正在编写的另一部著作中作更详细的阐述。
\end{enumerate}

\separation

在这篇导言中，我们不打算对代数几何中的“概形”观点作或详或略的描述，也不打算列出使我们不得不采用
这一观点的种种理由，特别是为何有必要系统地允许所研究的“簇”的局部环中存在幂零元。
这必然使有理映射的概念退居次位，而以正则映射或“态射”的概念居于首位。
本专著的目的恰在于系统地发展“概形”的语言，并且——我们希望——显示出它的必要性。
虽然这样做并不困难，我们在这里也不尝试对第一章所发展的概念作“直观的”引导。
希望预先了解本专著内容的读者，可以参阅 A. Grothendieck 在 1958 年 Edinburgh 国际数学家大会上的报告\cite{b7}，
以及同一作者的报告\cite{b8}。
J.-P. Serre 的论文\cite{b14}（记作 (FAC)）也可以看作介于代数几何的经典观点与概形观点之间的一种阐述；
就此而言，阅读该文可为阅读我们的《基础》作很好的准备。

\separation

为供参考，下面列出本专著预定的总体计划。这一计划以后仍可能修改，最后几章尤其如此：
\begin{enumerate}[label=第\Roman*章,align=left,leftmargin=5.3em,labelwidth=4.8em]
\item 概形的语言。
\item 若干类态射的初等整体研究。
\item 凝聚代数层的上同调及其应用。
\item 态射的局部研究。
\item 构造概形的初等方法。
\item 下降技术；构造概形的一般方法。
\item 群概形与主纤维空间。
\item 纤维空间的微分研究。
\item 基本群。
\item 留数与对偶性。
\item 相交理论、陈类与 Riemann—Roch 定理。
\item 阿贝尔概形与 Picard 概形。
\item Weil 上同调。
\end{enumerate}

原则上，各章都被视为开放的，随时可以在以后增加补充的节。
为了减少所采用的出版方式带来的不便，这些补充节将以单独的分册刊行。
当某章出版时，如果某个补充节已经列入计划或正在准备，即使出于轻重缓急的考虑，
其实际出版时间可能明显较晚，我们仍将在该章的目录中列出它。
为方便读者，我们在“第 0 章”中汇集一些关于交换代数、同调代数与层论的补充知识。
这些知识将在本专著各章中使用，虽然它们在不同程度上已为人所知，我们却未能找到便于引用的文献。
建议读者只在阅读本专著正文的过程中，并且只在对所引用的结果不够熟悉时，才查阅第 0 章。
我们还认为，这种读法会使本专著成为初学者熟悉交换代数和同调代数的一种良好途径：
倘若不伴随具体可感的应用，相当多的人会觉得这两门学科的学习枯燥，甚至令人沮丧。

\separation

在这篇导言中，即使只是对所阐述的概念和结果作简要的历史概述，也超出了我们的能力范围。
正文只列出我们认为对理解内容特别有用的文献，且只说明最重要结果的来源。
至少从形式上看，本书所论述的主题还相当新；这也解释了为什么我们很少引用十九世纪及二十世纪初
代数几何奠基者的工作——我们对这些工作的了解，仅来自转述。
不过，这里仍应略述那些最直接影响作者，并促成概形观点发展的著作。

首先必须提到 J.-P. Serre 的奠基性论文 (FAC)。对于不少被 A. Weil 经典著作
\textit{Foundations}\cite{b18} 的艰涩所吓退的年轻学习者而言，它是进入代数几何的门径；
本专著的一位作者也在其中。
该文第一次证明：“抽象”代数簇的“Zariski 拓扑”完全适合应用代数拓扑的某些技术，
尤其能够由此建立上同调理论。
此外，该文所给的代数簇定义，是最自然地适合扩展为我们在这里发展的概念的定义。\origfootnote{
据 J.-P. Serre 告知，应当指出：通过给定一个环层来定义簇的结构这一思想源于 H. Cartan，
他把这一思想作为其解析空间理论的出发点。当然，与代数几何中一样，在“解析几何”中，
也应当允许解析空间的局部环中出现幂零元。H. Grauert 最近已着手研究 H. Cartan 与 J.-P. Serre
定义的这种扩展\cite{b5}；我们有理由希望，不久将会出现一部在这一一般框架下系统阐述解析几何的著作。
显然，本专著所发展的概念和技术在解析几何中仍有意义，尽管应当预料到，后一理论中的技术困难将更为显著。
可以预见，代数几何由于其方法较为简单，能够为解析空间理论今后的发展提供某种形式上的模型。}

Serre 自己也曾指出：只要把域上的仿射代数换成任意交换环，仿射代数簇的上同调理论就可以毫无困难地转写过去。
因此，本专著的第一章、第二章以及第三章的头两节，就其主要内容而言，可以看作 (FAC) 以及同一作者
稍后的一篇论文\cite{b15} 中主要结果在这一扩大了的框架中的直接转写。
我们还从 C. Chevalley 的代数几何讨论班\cite{b1} 中获益良多；特别是系统地使用他引入的“可构造集”，
在概形理论中显示出很大的效用（参阅第四章）。
我们也借用了他从维数观点对态射所作的研究（第四章）；这一研究几乎无须作显著改动便可以转写到概形框架中。

另一方面，应当指出：Chevalley 引入的“局部环概形”概念，自然地适合用来扩展代数几何，
虽然它尚不具有我们在此希望赋予这一学科的全部灵活性和一般性。
关于这一概念与本书理论的关系，参见第一章第 8 节。
M. Nagata 已经在一系列论文\cite{b9} 中发展了这样的扩展，其中包含关于 Dedekind 环上的代数几何的许多特殊结果。\origfootnote{
在接近我们代数几何观点的工作中，应当提及 E. Kähler 的重要论文\cite{b22}，
以及 Chow 与 Igusa 最近的一篇短文\cite{b3}；后者在 Nagata—Chevalley 理论的框架内重新建立了 (FAC) 的若干结果，
并给出了一个 Künneth 公式。}

\separation

最后，不言而喻，一部关于代数几何的著作，尤其是一部关于其基础的著作，必然受到 O. Zariski 与 A. Weil
这样一些数学家的影响，哪怕这种影响是经由他人传递的。
特别是 Zariski 的全纯函数理论\cite{b20}，在借助上同调方法作适当调整，并补入一个存在定理
（第三章第 4、5 节）之后，与第六章阐述的下降技术一起，成为本专著所用的主要工具之一；
在我们看来，它也是代数几何现有的最强有力的工具之一。

这一理论所处的一般技术框架可以概述如下。第九章对基本群的研究将提供一个典型例子。
设有一个代数簇到另一个代数簇的固有态射（第二章）$f\colon X\to Y$，更一般地，也可以是一个概形到另一个概形的固有态射。
我们希望在一点 $y\in Y$ 的附近研究它，以解决一个关于 $y$ 的某邻域的问题 $P$。
研究按下列步骤进行：
\begin{enumerate}[label=\arabic*$^\circ$]
\item 可以假定 $Y$ 是仿射的，于是 $X$ 成为定义在 $Y$ 的仿射环 $A$ 上的概形；
甚至可以用 $y$ 的局部环来代替 $A$。在实际应用中，这一化归总是容易做到的（第五章），
它把问题归结为 $A$ 是局部环的情形。
\item 在 $A$ 是阿廷局部环时研究所考虑的问题。
为了使问题在不假定 $A$ 为整环时仍然确有意义，有时需要重新表述问题 $P$；
这样做往往反而能够使我们更好地理解问题，在这一阶段，这种理解具有“无穷小”的性质。
\item 形式概形理论（第三章第 3、4、5 节）使我们能够从阿廷环的情形过渡到完备局部环的情形。
\item 最后，若 $A$ 是任意局部环，则在 $X$ 上适当的概形中，考虑逼近某个给定“形式”截面的“多值截面”
（第四章），常常可以把对 $X$ 向 $A$ 的完备化作标量扩张后所得概形已经知道的结果，
转化为关于 $A$ 的一个足够简单的有限扩张——例如非分歧扩张——的类似结果。
\end{enumerate}

这个概述表明了系统研究阿廷环 $A$ 上的概形的重要性。
Serre 在表述局部类域论时所采用的观点，以及 Greenberg 最近的工作，似乎提示：
可以通过以函子方式给这样的概形 $X$ 配备一个定义在 $A$ 的剩余域 $k$ 上的概形 $X'$，来开展这种研究。
这里假定 $k$ 是完美域；在有利的情形下，$X'$ 的维数等于 $n\dim X$，其中 $n$ 为 $A$ 的长度。

至于 A. Weil 的影响，我们只须说明：写作本专著的主要动机之一，是需要发展必要的工具，
以便在所要求的一般性下表述“Weil 上同调”的定义，并着手证明\origfootnote{
为避免任何误解，须明确说明：在写这篇导言时，这项工作才刚刚开始，因此尚未得到 Weil 猜想的证明。}
为建立其著名的丢番图几何猜想\cite{b19} 所需的一切形式性质。
与此同样重要的动机，是希望为代数几何中的常用概念和方法找到自然的框架，
并使作者有机会理解这些概念与技术。

\separation

在结束这篇导言时，我们认为有必要预先提醒读者：与作者自己一样，读者大概也会经历一些困难，
才能习惯概形的语言，并确信几何直观所提示的通常构造，基本上都可以用唯一合理的方式转写成这种语言。
如同现代数学的许多领域一样，初始的直观与能够以所需的全部精确性和一般性来表达它的语言，
表面上正变得越来越远。
在这里，心理上的困难在于：我们必须把对集合而言熟悉的概念——笛卡儿积、群运算、环运算、模运算、
纤维空间、主齐性纤维空间等等——搬到一个已经颇不同于集合范畴的范畴的对象上，
也就是预概形范畴，或者给定预概形上的预概形范畴。
今后的数学家大概很难回避这种新的抽象努力；不过，归根结底，与我们的前辈为熟悉集合论而付出的努力相比，
它也许还相当轻微。

\separation

本书采用十进式层级编号给出引用。
例如在 III, 4.9.3 中，III 表示章，4 表示节，9 表示该节内的小节。
在同一章内部，省略章号。


% END chapters/introduction.tex